%% preliminaries.tex
%%

%% ==============
\chapter{Preliminaries}
\label{ch:preliminaries}
%% ==============

%This chapter should provide the foundations of the thesis.

%TODO maybe define a drawing too.
Before tackling the problem, we will introduce some concepts that are necessary 
for the understanding of this problem.
\begin{definition}[Graph Embedding]
    a Graph embedding of a graph $G$ on a  
\end{definition}

\begin{definition}[Level Assignement]
   A level assignement $l$ of a level graph $G$ is a mapping of vertices of $G$ to levels. 

    \[
        l: V(G) \rightarrow \mathbb{N}
    \]
\end{definition}
\begin{definition}[Proper Level Graph]
    a \textit{Proper Level Graph} $G$ is a graph with a \textit{level assignement} $l$ that 
    defines in which level every vertex is while respecting that for every edge $\left( u,v \right) \in E(G) \,\, |\,l(u) - l(v)\,| = 1$
\end{definition}

\begin{definition}[Planar Graph]
    a \textit{Planar Graph} $G$ is a graph that can have an embbeding on a 2-dimensional Euclidean 
    space $\mathcal{R}$ in a way that edges 
    only intersect with each other at their endpoints.
\end{definition}

\subsection{Equivalence relation and the 2-SAT formulation}%
\label{sub:Equivalence relation and the 2-SAT formulation}
Randerath et al.~\cite{randerath2001satisfiability} present some simple constraints that decide the 
existence of a level-planar drawing. The constraints are as follows: 
\begin{itemize}
    \item \textit{consistency:} $ uw \iff \neg wu   \quad \forall \, uw \in \mathcal{V} $
    \item \textit{transitivity:} $ \left( uw \wedge wx \right) \Rightarrow ux  \quad \forall \, uw, wx \in \mathcal{V} $ and with $ux \in \mathcal{V}$
    \item \textit{planarity:} $uw \iff vx \quad \forall uw, vx \in \mathcal{V} \text{ with }  \left( u,v \right), \left( w,x \right) \left( \in E \text{ and are independant} \right)$ 
\end{itemize}


\begin{definition}[Equivalence relation]
   An equivalence relation on a set $X$ is a binary relation $\sim$ on $X$ satisfying the following constraints: 
   \begin{itemize}
       \item \textit{reflexivity:} $a \sim a \, \forall a \in X$. 
       \item \textit{symmetry:} $a \sim b$ implies $b \sim a \quad \forall a, b \in X$.
       \item \textit{transitivity:} $a \sim b \wedge b \sim c $ implies $a \sim c \quad \forall a, b, c \in X$.
   \end{itemize}
\end{definition}

The planarity constraint can be formulated as an \textit{equivalence relation} over the set $\mathcal{V}$ according to Randerath et
al.~\cite{randerath2001satisfiability}.

\begin{definition}[Orbit relation]
    \label{def:orbit}
\ \\ 
     Let's consider a proper level graph $G = \left( V,E \right)$, and a boolean set 
    $\mathcal{V} = \left\{ uw \; | \; u,w \in V,\, l\left( u \right)=l\left( w \right)   \right\}$. \\
    The orbit relation is an equivalence relation on the boolean set $\mathcal{V}$ with a binary 
    relation $\sim$ on $\mathcal{V}$ satisfying the following constraints:  
    \begin{itemize}
        \item \textit{reflexivity:} $ uw \sim uw   \quad \forall \, uw \in \mathcal{V} $
        \item \textit{symmetry:} $ uw \sim vx $ implies $ vx \sim uw \quad \forall \, uw, vx \in \mathcal{V} $
        \item \textit{transitivity:} $ \left(uw \sim zy \right) \wedge  \left( zy \sim vx \right) $ implies $ uw \sim vx \quad \forall \, uw, zy, vx \in \mathcal{V} $
    \end{itemize}
\end{definition}

The interperation of the orbit of a variable $uw \in \mathcal{V}$ is that for each boolean variable $vx$ 
of $\mathcal{V}$ in the orbit of $uw$ we have: $\left( u < w \right) \iff \left( v < x \right)$.\cite{randerath2001satisfiability}

\begin{theorem}
    \label{th:randert6}
    A proper level graph $G$ is level planar iff there exists no $uw \in \mathcal{V}$ such that $wu$ in orbit of $uw$. 
\end{theorem}
\begin{proof}
    Refer to Theorem 6 from~\cite{randerath2001satisfiability}.
\end{proof}

%which means for a level planar graph, there is an other constraint that (paper) named consistency constraint that is also

by considering the orbit which represents the \textit{planarity} constraint and by considering the \textit{consistency} constraint, 
that is necessary to prove the existence of a level planar based on the Theorem~\ref{th:randert6}. Randert et al~\cite{randerath2001satisfiability} 
proves that the transitivity clause is not needed for the satisfiability of this system of constraints and it is equivalent to 
the satisfiability of a reduced constraint system which omits the transitivity clause. 

The reduced constraint system can be expressed in terms of 2-SAT, which can be solved efficiently with a quadratic-time 
algorithm for testing level planarity~\cite{krom1967decision}. Randert et al. also prove that a solution can be made transitive 
without invalidating the other constraints in Theorem~\ref{th:randert6}.

\subsection{Hanani-Tutte drawing and Embedding}%
\label{sub:Hanani-Tutte drawing and Embedding}

\begin{theorem}
    \label{th:strongHT}
   If $G$ has an (x-) monotone drawing with every pair of independant edges crosses evenly then $G$ has (x-) monotone
   embedding with the same vertex locations (over x). 
\end{theorem}
\begin{proof}
    Refer to theorem 3.1 from~\cite{fulek2012hanani} 
\end{proof}

\subsection{Equivalence classes of the orbits and Hanani-Tutte drawing }%
\label{sub:eqcl&htd}
In this section we want to link the equivalence classes with the Hanani-Tutte drawing induced from the truth assignement 
of the 2-SAT formulation that is based on equivalence classes of the orbits.  

\begin{definition}[Equivalence class]
    For a given set $X$ and an equivalent relation $\sim$ in $X$, for an element $x$ of $X$, the 
    equivalence class $[x]_\sim = \left\{ y \sim x, y \in X\right\}$. 

    The set of all equivalence classes of $X$ with a relation $\sim$ is represented with $X/\sim$ 
\end{definition}

\begin{definition}[Orbit class]
    For a boolean variable $uw \in \mathcal{V}$, the orbit class $[uw]_\mathcal{\sim}$ is 
    the set of all elements $vx \in \mathcal{V}$ with $vx \sim uw$. 

    The set of all orbit classes of $\mathcal{V}$ with the relation $\sim$ is $\mathcal{V}/\sim$.

\end{definition}

\begin{lemma}
    $\forall uw \in \mathcal{V} \  \forall vx \in [uw]_{\sim} : \neg vx \in [\neg uw]_{\sim}$  
\end{lemma}
\begin{proof}
    by the consistency clause, $\forall uw \in \mathcal{V} : uw \iff \neg wu$ and  $\forall vx \in [uw]_{\sim} vx \iff uw$ 
    therefore \[\forall uw \in \mathcal{V} \; \forall vx \in [uw]_{\sim} vx \iff \neg wu\]
    thus \[\forall uw \in \mathcal{V} \; \forall vx \in [uw]_{\sim} \neg vx \iff wu \iff \neg uw \] 
    and that's why \[\forall uw \in \mathcal{V} \; \forall vx \in [uw]_{\sim} \neg vx \in [\neg uw]_{\sim} \] 
\end{proof}

Each assignment of the equivalence classes of the orbit is represented in the hanani-tutte drawing $\Gamma^*$ of $G^*$ in the 
level $l$ where the single vertex $v'$ of origin $v$ in $G$ is one of the pair of vertices in the equivalence class and the other vertex $w$
of the pair is the edge $e'$ passing through the level $l$ with an endpoint of origin $w$ in the graph $G$. 
If for example we assign one to $[vw]_{\sim}$, then the vertex $v'$ lays before $e'$ on the level $l$ in the drawing $\Gamma^*$. 

\begin{definition}[Horizental Gang]
    a \textit{Horizental Gang} is a set of distinct vertices that belong to the same level of the level graph $G$ and that are the distinct 
   endpoints of adjacent edges.
\end{definition}

\begin{definition}[Base Orbit Class]
    a base orbit class $\mathcal{B}\left( [uv]_\sim \right)$ is the merge of $[uv]_\sim$ and $[vu]_\sim$. 
    The orbit classes $[uv]_\sim$ and $[vu]_\sim$ are the directions of the base orbit class $\mathcal{B}\left( [uv]_\sim \right)$.
    
\end{definition}
\begin{lemma}
    \label{lm:cyclerel}
   Let a level planar graph $G=\left( V,E \right)$ with a level assignment $l$ and let the drawing $\Gamma^*$ be the induced 
   Hanani-Tutte drawing of the graph $G^*$ induced from the orbit classes of $G$. An odd crossing in $\Gamma^*$ is 
   equivalent to a cyclic relation between a triple of members of a horizental gang $V_{HG}$, in other words
   a pair or triple of base orbit classes of members of the horizental gang can have directions that are cyclic. 
\end{lemma}
\begin{proof}
   TODO proof 
\end{proof}

The algorithm that we designed might merges some of the orbit classes, we want to seperate the orbit classes before 
the execution of the algorithm from any other state of the $\mathcal{V}/\sim$, for this reason we introduce this definition. 

\begin{definition}[first-gen orbit class]
    A \textit{first-gen} orbit class is an orbit class purely based on the planarity constraint, wich is not contaminated 
    with other elements of $\mathcal{V}$ that are not par of it from the planarity constraint.
\end{definition}


